\chapter{X-ray, Far-Infrared, Submillimeter, and Radio Diagnostics}
\label{chap:multiwavelength}

\section{One gas cycle, several observational weightings}

No band inventories every phase of cosmic gas. Optical and ultraviolet
spectra are highly sensitive to $10^4$ K ionized gas and to absorption through
low-density material, but dust, cold molecular reservoirs, and hot plasma can
remain inconspicuous. X-ray lines expose highly ionized metals; far-infrared
fine-structure lines carry cooling from neutral and weakly ionized gas;
millimeter rotational lines make cold molecular gas observable; and centimeter
data trace atomic hydrogen, free--free emission, synchrotron radiation, and
magnetized plasma. These are complementary density, temperature, abundance,
opacity, and geometry weightings, not interchangeable mass estimators.

The common discipline is forward modeling. A source spectrum is altered by
radiative transfer and cosmological redshift, sampled through a point-spread
function or synthesized beam, and redistributed among detector channels.
Continuum, aperture, response, and background therefore belong in the
definition of every reported line measurement.

\section{Thermal X-ray plasma}

For a low-density, optically thin collisionally excited line,
\begin{equation}
 L_{ul}=\int n_{\rm e}n_{\rm H} A_X f_{X^i}(T,n_{\rm e}t)
 \epsilon_{ul}(T,n_{\rm e})\,{\rm d}V ,
 \label{eq:xray-line}
\end{equation}
where $A_X$ is elemental abundance, $f_{X^i}$ the ionic fraction, and
$\epsilon_{ul}$ the emissivity coefficient. The $n_{\rm e}n_{\rm H}$
weighting favors dense structure, so a uniform-density conversion from
luminosity to mass is biased when gas is clumpy.

In collisional ionization equilibrium (CIE), ionization balances
recombination and ion fractions depend mainly on temperature. A recently
shocked plasma may be underionized; rapidly cooling or expanding gas may be
overionized and recombining. Nonequilibrium calculations introduce an
ionization age such as $n_{\rm e}t$. Photoionized plasma instead responds to
an external spectral energy distribution and ionization parameter. One ion
can occur in more than one regime, so a temperature assigned from a line name
alone is model dependent \citep{Draine2011,SutherlandDopita1993,AtomDB2025}.

H-like Ly$\alpha$ lines, He-like triplets, the Fe L complex, and Fe K lines
provide complementary constraints. For a He-like triplet,
\begin{equation}
 G=\frac{f+i}{r},\qquad R=\frac{f}{i},
\end{equation}
where $r$, $i$, and $f$ denote resonance, intercombination, and forbidden
components. These ratios respond to excitation mechanism, temperature,
density, and ultraviolet pumping, but blends and radiative transfer must be
included.

A multiphase source can be represented by a differential emission measure,
\begin{equation}
 L_E=\int \Lambda_E(T,Z)
 \frac{{\rm d}EM}{{\rm d}\log T}\,{\rm d}\log T,\qquad
 EM=\int n_{\rm e}n_{\rm H}\,{\rm d}V .
\end{equation}
Fitting a one-temperature spectrum to such gas can bias both continuum and
abundance. The adopted solar abundance scale must be named
\citep{Asplund2021}. Likewise, the cooling relation
$\mathcal{L}=n_{\rm e}n_{\rm H}\Lambda(T,Z)$ hides dependencies on ionization
state, abundance pattern, and radiation field.

\subsection{Absorption, fluorescence, and the detector response}

Photoelectric absorption may be written schematically
\begin{equation}
 F_E^{\rm obs}=F_E^{\rm int}
 \exp[-N_{\rm H}\sigma_{\rm eff}(E,Z,\xi)].
\end{equation}
The fitted equivalent hydrogen column inherits assumptions about abundances,
cross-sections, ionization, and covering. Fe K$\alpha$ near $6.4$ keV can
arise in distant obscuring matter or disk reflection. Near a black hole,
Doppler shifts, gravitational redshift, and light bending broaden the full
reflection spectrum; an isolated Gaussian is a detection model, not a
complete physical interpretation \citep{GeorgeFabian1991,Reynolds2014}.

Counts predicted in channel $j$ are
\begin{equation}
 \mu_j=t_{\rm exp}\int R_j(E)A_{\rm eff}(E)S(E)\,{\rm d}E+B_j ,
\end{equation}
so photon models must be folded through redistribution and effective-area
responses. Sparse counts call for Poisson likelihoods rather than forced
Gaussian errors. A free search in line energy also requires a trials
correction. Weak O~VII or O~VIII absorption claims in the WHIM are especially
sensitive to continuum choice, wavelength calibration, fixed-pattern
residuals, Galactic foregrounds, and the full searched redshift path
\citep{Bregman2007,Nicastro2018}.

Microcalorimeters resolve velocity information without dispersion. XRISM
Resolve provides few-eV-class spectroscopy over much of its operating band
\citep{XRISMPOG2026}. For approximately Gaussian components,
\begin{equation}
 \sigma_v^2=\frac{k_{\rm B}T}{m_{\rm ion}}+
 \sigma_{\rm turb}^2+\sigma_{\rm bulk,unresolved}^2 .
\end{equation}
Ions of different mass can help separate thermal and nonthermal widths only
if they trace the same volume.

\section{Far-infrared fine-structure cooling}

Low-lying fine-structure transitions can be excited in gas too cool for
optical forbidden lines. Central examples are [C~II] $158\,\mu{\rm m}$,
[O~I] $63$ and $145\,\mu{\rm m}$, [N~II] $122$ and $205\,\mu{\rm m}$,
[O~III] $52$ and $88\,\mu{\rm m}$, and [C~I] $370$ and $609\,\mu{\rm m}$.
For an optically thin two-level approximation,
\begin{equation}
 n_{\rm crit}=\frac{A_{ul}}{\sum_c q_{ul}^{(c)}}.
\end{equation}
A critical density is incomplete unless collider and temperature are stated.
Electrons, H, and H$_2$ have different rate coefficients, and radiative
trapping changes the effective excitation requirement.

[C~II] is not phase unique: C$^+$ occurs in ionized gas, diffuse neutral gas,
and photodissociation-region (PDR) surfaces. [N~II] helps estimate the ionized
contribution, while [O~I] weights warmer, denser neutral material but may be
optically thick or self-absorbed. [O~III] selects harder radiation. A PDR-grid
solution is conditional on geometry, abundance, column, filling factor, and
the partition of [C~II] among phases \citep{TielensHollenbach1985,
HollenbachTielens1999}.

The observed decline of [C~II]/far-infrared luminosity in many compact
luminous systems is called the [C~II] deficit. High ionization parameter,
grain charging, dust-bounded regions, saturation of the upper level, opacity,
and an AGN-heated continuum may all contribute. The name describes a ratio;
it does not select one cause.

\section{Molecular rotational lines}

Cold H$_2$ has no permanent dipole moment, so CO is widely used as a proxy.
The velocity-integrated CO flux gives the brightness-temperature luminosity
\begin{equation}
 L'_{\rm CO}=3.25\times10^7
 (S_{\rm CO}\Delta v)\nu_{\rm obs}^{-2}D_{\rm L}^{2}(1+z)^{-3}
 \quad {\rm K\,km\,s^{-1}\,pc^2},
\end{equation}
for flux in Jy km s$^{-1}$, frequency in GHz, and distance in Mpc. Molecular
mass is then
\begin{equation}
 M_{\rm mol}=\alpha_{\rm CO}L'_{\rm CO(1-0)}.
\end{equation}
The factor $\alpha_{\rm CO}$ depends on metallicity, surface density,
temperature, dynamics, cloud structure, and the helium convention; it is a
physical prior rather than a unit conversion \citep{Bolatto2013}.

A high-$J$ line additionally requires an excitation ratio
$r_{J1}=L'_{\rm CO[J-(J-1)]}/L'_{\rm CO(1-0)}$. A CO spectral-line energy
distribution constrains a distribution of density and temperature and rarely
has a unique one-component interpretation. Dense-gas tracers such as HCN,
HCO$^+$, and CS are also affected by opacity, trapping, chemistry, and
infrared pumping \citep{CarilliWalter2013}.

At high redshift, $T_{\rm CMB}=2.725(1+z)$ K both excites low levels and
provides a background. Ignoring it can hide cold gas whose excitation
temperature lies only slightly above the background. Main-isotope CO is
often optically thick; isotopologues constrain opacity but introduce isotope
chemistry and selective photodissociation.

\section{H~I at 21 cm}

The hyperfine line has rest frequency $1420.40575177$ MHz. A uniform layer
obeys
\begin{equation}
 T_{\rm B}(v)=T_{\rm s}(1-e^{-\tau})+T_{\rm bg}e^{-\tau}.
\end{equation}
For optically thin foreground-subtracted emission,
\begin{equation}
 N_{\rm HI}=1.823\times10^{18}\int T_{\rm B}(v)\,{\rm d}v
 \quad {\rm cm^{-2}},
\end{equation}
with $T_{\rm B}$ in K and velocity in km s$^{-1}$. In absorption,
\begin{equation}
 N_{\rm HI}=1.823\times10^{18}\frac{T_{\rm s}}{f_{\rm c}}
 \int\tau_{\rm obs}(v)\,{\rm d}v .
\end{equation}
Thus absorption alone constrains a combination of spin temperature and
covering factor \citep{Chengalur2013}. For unresolved optically thin emission,
\begin{equation}
 \frac{M_{\rm HI}}{M_\odot}=2.356\times10^5
 \frac{D_{\rm L}^2}{1+z}
 \left(\frac{\int S_\nu{\rm d}v}{\mathrm{Jy\,km\,s^{-1}}}\right),
\end{equation}
with $D_{\rm L}$ in Mpc. Self-opacity, beam dilution, stray radiation, and
resolved-out diffuse structure violate simple conversions.

\section{Radio continuum, magnetic fields, and interferometers}

Optically thin free--free radiation is tied to the ionizing-photon rate,
whereas synchrotron emission depends on cosmic-ray electrons, magnetic
fields, losses, and transport. Their spectral decomposition is approximate
when regions are mixed or when low-frequency absorption and curvature matter
\citep{CondonRansom2016,KennicuttEvans2012}.

For a simple foreground screen,
\begin{equation}
 \chi(\lambda)=\chi_0+{\rm RM}\lambda^2,\qquad
 {\rm RM}=0.812\int n_{\rm e}B_\parallel\,{\rm d}l
 \quad {\rm rad\,m^{-2}},
\end{equation}
with $n_{\rm e}$ in cm$^{-3}$, $B_\parallel$ in $\mu$G, and path in pc.
Internal rotation, multiple components, and beam depolarization require
Faraday-depth modeling rather than one slope.

An interferometer samples Fourier components:
\begin{equation}
 V(u,v)=\int A(l,m)I(l,m)e^{-2\pi i(ul+vm)}\,{\rm d}l\,{\rm d}m .
\end{equation}
The beam, weighting, flagging, and deconvolution are part of the measurement.
Long baselines sharpen resolution but do not ensure surface-brightness
sensitivity; absent short spacings suppress large-scale flux. A reproducible
line cube records channel response, synthesized beam, primary-beam correction,
continuum subtraction, mask construction, and correlated noise.

\section{A cross-band audit}

The strongest combinations break different degeneracies: H$\alpha$ plus
free--free emission tests extinction; [N~II] partitions [C~II]; CO ladders
constrain excitation while dust supplies an independent mass prior; and H~I
emission plus absorption constrains spin temperature. Comparisons require
matched aperture, beam, velocity interval, and calibration convention.

Always preserve the inference chain
\[
 \begin{gathered}
 \hbox{counts or visibilities}\rightarrow\hbox{line flux and profile}
 \rightarrow\hbox{excitation and opacity}\\
 \rightarrow\hbox{column or mass}\rightarrow\hbox{physical claim}.
 \end{gathered}
\]
A bright line need not trace the dominant mass, a nondetection is not a
phase-independent zero, critical density is not a hard boundary, and a
higher-resolution interferometric image may recover less total flux. Every
arrow in the chain adds assumptions; an atlas should display them beside the
diagnostic.
